\section{Exact saturation and sharp growth deficits}\label{sec:saturation}

The following statement applies to any finite downset of compositions,
not only to a principal ideal or a monotone host. Let $f_k$ count its
members of weight $k$.

\begin{theorem}[Saturation rigidity at every rank]
Let $k\ge1$ and assume $f_k>0$. If $f_{k+1}=2f_k$, then the downset
contains every composition of weights $k$ and $k+1$. In particular,
$f_k=2^{k-1}$ and $f_{k+1}=2^k$.
\end{theorem}

\begin{proof}
Decrement the last part of a word of weight $k+1$, deleting it if it is
one. Every word $v$ of weight $k$ has at most two preimages: append a
one, or increase its last part. Equality of the total counts forces
both preimages of every feasible $v$ to be feasible.

Now remove one unit from the first part of either preimage, deleting
that part if necessary. Downward closure keeps the resulting word
feasible. Encode a composition of weight $k$ by its $k-1$ cut bits.
These two operations delete the first cut bit and append, respectively,
a one or a zero. Starting from any feasible cut string, $k-1$ successive
operations, with the appended bits prescribed, produce any desired cut
string of length $k-1$. Thus every composition of weight $k$ is feasible.
For $k=1$ the cut string is empty and the same conclusion is immediate.
The two already forced extensions then give every composition of
weight $k+1$.
\end{proof}

\begin{corollary}[The exact support of the growth deficit]
For a nonempty host $w$ of weight $N$, let $\kappa(w)$ be the smallest
weight of a composition not below $w$. Put
\[
 E_w=1-z+(2z-1)F_w.
\]
Then the positive coefficient support of $E_w$ is precisely
$[\kappa(w),N+1]\cap\mathbb Z$. Consequently
\[
 (F_w)_{n+1}\le2(F_w)_n-1
 \quad\text{for }\kappa(w)-1\le n\le N.
\]
Moreover $2\le\kappa(w)\le\min\{\len w+1,\max w+1\}$.
\end{corollary}

\begin{proof}
All ranks through $N$ are nonempty, and the two-preimage map gives
$f_{j}\le2f_{j-1}$ for $j\ge2$. Saturation rigidity identifies equality
exactly with the presence of every composition of rank $j$. Fullness
of a rank implies fullness of every smaller rank, by downward closure.
The coefficients at degrees zero and one cancel; beyond $N+1$ they
vanish. The words $1^{\len w+1}$ and $(\max w+1)$ give the upper bounds
on $\kappa$, and every nonempty host contains $(1)$.
\end{proof}

\begin{lemma}[Threshold decomposition at a maximum]
Set $\kappa(\varnothing)=1$. For an arbitrary composition
$w=(X,M,Y)$ with $M\ge\max(X,Y)$,
\[
 \kappa(w)=\min\{M+1,\kappa(X)+\kappa(Y)\}.
\]
\end{lemma}
\begin{proof}
The singleton $(M+1)$ is absent. If $u$ and $v$ are missing compositions
of minimum weights for $X$ and $Y$, respectively, then $uv$ is absent
from $XMY$: an embedding of $u$ cannot finish inside $X$, leaving
an embedding of $v$ entirely inside $Y$, a contradiction.

Conversely, take any composition of weight
$k<\min\{M+1,\kappa(X)+\kappa(Y)\}$. If $k<\kappa(X)$ it embeds
in $X$. Otherwise choose its first part for which the cumulative weight
reaches $\kappa(X)$. Its preceding parts have weight below $\kappa(X)$
and embed in $X$; its following parts have weight at most
$k-\kappa(X)<\kappa(Y)$ and embed in $Y$. The chosen part is at most
$k\le M$ and embeds in the central maximum.
\end{proof}

This decomposition uses the same maximum-layer universality mechanism
as Albert, Engen, Pantone and Vatter~\cite{AlbertEtAl}, proof of Theorem~1.
We use it to identify the parameter in the growth
bound; the universality construction is not claimed as new here.
It also gives a linear-time procedure on a maximum Cartesian tree of
the host, and a minimum-weight missing composition as a witness.

\begin{proposition}[Universality threshold for a monotone arm]
If $Y=(t_1\le\cdots\le t_l)$ is nonempty, the largest weight at which
all compositions occur is
\[
 \kappa(Y)-1=\min\left\{l,\ \min_{1\le i\le l}(i-1+t_{l-i+1})\right\}.
\]
\end{proposition}

\begin{proof}
Fix a weight $k$. Necessity of $l\ge k$ follows from $1^k$. For every
$1\le i\le k$, the word $(k-i+1,1^{i-1})$ requires
$t_{l-i+1}\ge k-i+1$. These inequalities are also sufficient: a word
of weight $k$ and length $i$ has every part at most $k-i+1$, and hence
embeds in the last $i$ positions. The inequalities for $i>k$ are
automatic, so the displayed minimum is equivalent to all the necessary
and sufficient conditions.
\end{proof}

The exact support statement identifies the first strict-growth rank
without a monotonicity hypothesis on the host. We next quantify the
deficit at every subsequent rank.
Enumeration and equivalence of single-pattern avoidance classes in
this order are studied by Menon and Singh~\cite{MenonSingh}. The
bounds below apply to every downset inside such an avoidance class,
not only to the full avoidance class.

\begin{proposition}[A quantitative saturation bound]
For every finite composition downset and every $k\ge2$, set
$s=f_k$ and $N=2^{k-1}$. Then
\[
 (k-1)N(2f_k-f_{k+1})\ \ge\ 2s(N-s).
\]
In particular, if $s>0$, then
\[
 \frac{f_{k+1}}{f_k}
 \le 2-\frac{2}{k-1}\left(1-\frac{f_k}{2^{k-1}}\right).
\]
\end{proposition}

\begin{proof}
Put $n=k-1$. Encode feasible words of weight $k$ as a set $S$ of binary
strings of length $n$. A feasible word of weight $k+1$ determines an
edge $x\longrightarrow x_2\cdots x_n b$ between two members of $S$:
the endpoints are obtained by removing the last and the first unit of
the word. The edge determines the word uniquely. Consequently
$2s-f_{k+1}$ is at least the number $e$ of directed edges from $S$ to
its complement in the binary shift graph.

There is a unique directed path of length $n$ from any initial string
$x$ to any terminal string $y$, obtained by appending the bits of $y$.
The $s(N-s)$ paths with $x\in S$ and $y\notin S$ each traverse an edge
from $S$ to its complement. Fix an edge and a position $i\in\{1,\ldots,n\}$
at which that edge is traversed. Without restrictions on the initial
and terminal strings, the number of such length-$n$ paths is
$2^{i-1}2^{n-i}=N/2$. A union bound over the $e$ boundary edges and
the $n$ possible positions gives $s(N-s)\le neN/2$.
Combine this with $e\le2s-f_{k+1}$.
\end{proof}

\begin{theorem}[Sharp uniform deficit from a missing composition]
Suppose a composition of weight $\kappa\ge2$ is absent from a
composition downset. Then, for every $k\ge\kappa-1$,
\[
 2f_k-f_{k+1}\ \ge\ \frac{f_k}{2^{\kappa-1}-1}.
\]
The assertion applies to infinite downsets as well, since each weight
contains only finitely many compositions. The constant is best possible
for every $\kappa$.
\end{theorem}

\begin{proof}
At $k=\kappa-1$, saturation rigidity and the missing word imply
$2f_k-f_{k+1}\ge1$ whenever $f_k>0$, while
$f_k\le2^{\kappa-2}\le2^{\kappa-1}-1$. The zero case is immediate.
It remains to consider $k\ge\kappa$.
Use the preceding binary shift graph on strings of length $n=k-1$,
and put $h=\kappa-1\le n$. Append the $h$ cut bits of a fixed missing
composition $v$. The resulting composition contains $v$: its last
$\kappa$ units have the prescribed internal cuts, and the first part
of $v$ may only be enlarged by preceding units. Hence the terminal
string is outside $S$. Starting from each of the $s=f_k$ members of
$S$ gives $s$ paths of length $h$, each meeting an outgoing boundary
edge. At a fixed position $i$, a fixed edge can be traversed by at
most $2^{i-1}$ of these paths, because the appended block
is fixed. Summing over the $h$ positions and the
$e$ outgoing edges gives $s\le e\sum_{i=1}^h2^{i-1}=e(2^h-1)$. Now use
$e\le2f_k-f_{k+1}$.

For sharpness, take the downset of all compositions with parts at most
$\kappa-1$, and compare weights $k=\kappa$ and $k+1$.
Exactly one composition is excluded at weight $\kappa$, whereas
exactly three are excluded at weight $\kappa+1$, namely
$(\kappa+1)$, $(\kappa,1)$ and $(1,\kappa)$. Thus
$f_\kappa=2^{\kappa-1}-1$ and
$2f_\kappa-f_{\kappa+1}=1$. A finite principal ideal gives the same
example: use the rectangular host $(\kappa-1)^{\kappa+1}$.
\end{proof}
